% !TeX root = ../papers/nc-p3.tex
\section{有限な低階数軌跡を持つ可約な反対八次曲面の因子}
\label{sec:reducible-octic-factors}
\begin{theorem}\label{thm:reducible-octic-factors}
四辺有限の標準有限Frobenius窓で入力 $X$ は被約、全階数反対平方根 $N$ の階数2以下の軌跡は有限とする。
$d=\det N$ は平方因子を持たず、次数4以上の既約因子を持つ。
各 $p\in X$ で重複度はちょうど3、接錐の射影三次曲線は被約である。
因子 $S$ の $X$ における重複度の和 $m(S)$ は、次数 $1,2,3$ に対しそれぞれ $7,14,22$ 以下である。
\end{theorem}
\begin{proof}[例外平面と行列式の重複]
定理\ref{thm:double-point-opposite-reconstruction}により六次余因子は $\mathcal I_X^2$ を生成する。
$B=\operatorname{Bl}_X\mathbb P^3$ で超平面を $H$、例外平面の和を $E$ とし、$P=3H-E$ と置く。
$2P$ は基点を持たず、$P$ はnefである。各点で単元成分を消した対称三行の一次部を $A(u)$ とする。
六つの二次小行列式は例外平面上の全二次式を張るから階数は至るところ2以上であり、
定理\ref{thm:octic-kernel-image-degree}の局所証明より $\det A$ は零多項式でない。
従って $d$ の初期形式は単元倍の $\det A$ である。その三次曲線 $C$ 上で余核は線束であり
\[
 0\longrightarrow\mathcal O(-2)^3\xrightarrow{A}\mathcal O(-1)^3
 \longrightarrow i_*\theta\longrightarrow0
\]
を得る。実際各点で単元二行を消去できる。$H^0(\theta)=0$ で、$\mathcal O(-3)$ による双対と
$\omega_C=\mathcal O_C$ により $\theta\simeq\theta^\vee$ である。非被約三次ならその被約部は
直線または交わる二直線で二捩れを持たない。冪零イデアル上の単元群も標数零で対数により
加法群となり、定数単元は持ち上がるのでPicard群の制限核に二捩れはない。従って $C$ は被約である。
同じ理由で既約尖点三次も除かれる。実際その正規化は射影直線で導手の単元商は加法群である。

重複因子があればその因子はどの $p$ も通れず、二乗で割って得る次数 $t\le6$ の式がなお全点を三重に通る。
しかし有効因子 $tH-3E$ について $P^2(tH-3E)=9t-60\ge0$ が必要で、これは矛盾する。
\end{proof}
\begin{proof}[各因子における三次式と核座標]
次数 $e$ の既約因子 $S$ につき $n=|X\cap S|$, $m=m(S)$ とする。
その真変換上で余因子射はVeroneseを経る。従ってこの真変換の滑らかな射影解消 $T$ 上で
核座標の生成する線束 $L$ は $2L=2P|_T$ を満たし、
\[
                  L^2=P^2[S']=9e-m
\]
となる。正なら核像は曲面で $h^0(L)\ge3$ である。$T$ が有理曲面ならPicard群に捩れがなく $L=P|_T$。
この切断は $3H$ の切断として各印の逆像で消える。$S$ が正規なら $H^0(T,3H)=H^0(S,\mathcal O_S(3))$ で、
この空間は周囲の三次式により尽くされ、$n$ 個の値条件は独立である（定理\ref{thm:finite-cubic-separators}）。
この議論は有理性を解消のPicard群に二捩れがないという仮定に置き換えても成立する。

平面では定理\ref{thm:finite-plane-section-seven}より $m=n\le7$。
任意の非零二次式 $q$ に対し $q\ne0$ である印の数 $r$ は6以上である。
実際 $r>0$ は三次評価の単射性による。隣の辺の $(2,2)$ 完全跡を $q$ 倍すると、
この $r$ 点上の完全跡が隣の一次座標空間の平方を消す。従って一次座標の評価像の次元は $r/2$ 以下。
$r=1,2,3$ では点への閉埋込み、$r=4,5$ では直線切断長3以下に反する。
よって二次因子は高々14点を含む。既約二次曲面の重複度は一つの頂点を除き1なので $m\le15$。
もし $n=14$ なら有理正規曲面の解消上で $L^2\ge3$ にもかかわらず
$h^0(L)\le16-14=2$ となる。故に $n\le13$, $m\le14$ である。
\end{proof}
\begin{proof}[三次因子と特異母線]
三次因子では $n\le19$。孤立した二重点しか持たない三次曲面は特異点が高々四個である。
三点が共線なら全偏導二次式が直線を含み、四点が共面なら三点は共線でないから六本の割線全てを
平面三次が含んでしまう。四非共面点を座標点に選ぶと式は $\sum_i c_i\prod_{j\ne i}x_j$ であり、
既約性から全 $c_i\ne0$、これらの座標点以外は滑らかである。
従って $m\le23$。等号なら $n=19$ で、二重点からの射影で有理、正規なので
$L^2=4$ と $h^0(L)\le19-19=0$ が矛盾する。三重点を持てば平面三次上の錐であり、
底が滑らかなら $m\le21$ である。

特異曲線を持つ場合、その曲線は一本の直線で他に特異点はない。
実際二特異点の割線は三次曲面に含まれる。非直線な平面曲線なら割線が平面を覆い、
非平面曲線の割線は三次元を覆う（一般点からの射影、または割線の次元定理）。
特異直線とその外の特異点からも平面が生じる。三重点がなければ直線切断長から $m\le19+3=22$。
残る $m\ge23$ は特異平面三次上の錐で $m=23$, $n=19$、頂点と特異母線上の他の二点が印の場合だけである。
頂点の接三次に関する第一段より底は節点三次である。正規化した底の上の錐の解消は
$\mathbb F_3$ で、超平面は $H=C_0+3f$（$C_0^2=-3$）となる。
共通解消を取って $L=P$ の切断を見ると、これは $\mathbb F_3$ 上の $3H$ の切断であり
各印の全逆像で消える。従って19個の印の上で一つずつ選んだ値条件に加え、
二つの非頂点の特異な印それぞれで両枝の値の差も零である。

これらは少なくとも21個の独立条件を与える。周囲の三次式は初めの19値に全射し、両枝の差を零にする。
一方、$\mathbb F_3$ 上では切断を $\sum_{j=0}^3 w^{3-j} a_j$,
$a_j\in H^0(\mathbb P^1,\mathcal O(3j))$ と書ける。$w$ は頂点以外の母線のアフィン座標であり、
節点の二つの持ち上げで底の超平面束の自明化を貼合せに合わせる。
$j=1,2,3$ で二点間の係数の差を独立に指定でき、従って差は $w$ の任意の二次以下の式となる。
二印の $w$ は異なるので二つの値に全射する。ゆえに $h^0(L)\le(1+4+7+10)-21=1$ で $L^2=4$ に反する。
以上で $m\le22$ となった。全因子で $m$ の和は60であるが、全て三次以下なら三次因子の個数 $a\le2$ に対し
和は高々 $7(8-3a)+22a=56+a<60$。これで次数の結論を得る。
\end{proof}

\path{computations/verify_reducible_octic_factors.py} は吹上げの交叉数、節点錐の二枝の評価および因子次数分割の上界を検算する。
